\chapter{Introducción}
El álgebra de Boole, originalmente introducida por George Boole en su obra de 1854 \textit{The Laws of Thought}, fue posteriormente axiomatizada de forma rigurosa. En 1904, el matemático estadounidense Edward Vermilye Huntington presentó un conjunto de postulados independientes que definen formalmente una estructura de álgebra de Boole. 

Para evitar cualquier confusión conceptual con la aritmética tradicional, en esta fase inicial emplearemos la signatura propia de la teoría de retículos, utilizando los símbolos $\vee$ (supremo o join) y $\wedge$ (ínfimo o meet), junto a los elementos constantes $\bot$ (mínimo) y $\top$ (máximo). El complemento se denotará con el símbolo clásico de la negación lógica $\neg$. Más adelante, y por conveniencia práctica, transitaremos hacia la notación clásica de sistemas digitales ($+$, $\cdot$, $0$, $1$).

\section{Pre-Axiomas de la Estructura}
Antes de enunciar los postulados, debemos definir rigurosamente sobre qué elementos y operaciones estamos trabajando. 

Partimos de un ente matemático $\mathbb{B}$.

\begin{preaxioma}{Estructura de Conjunto - EsConjunto($\mathbb{B}$)}{esconj}
Se requiere que $\mathbb{B}$ sea un conjunto.
\end{preaxioma}

\begin{preaxioma}{Elementos Constantes - Constantes}{constantes}
Este conjunto ha de cumplir que tiene dos elementos que llamaremos constantes, tales que $\bot \in \mathbb{B}$ y $\top \in \mathbb{B}$. En principio, no asumimos nada sobre la igualdad o desigualdad de estas constantes.
\end{preaxioma}

Además, vamos a definir dos operaciones binarias internas que denotaremos por $\vee$ y $\wedge$. Estas deben satisfacer rigurosamente la definición de función:

\begin{preaxioma}{Operación Binaria Interna $\vee$ - OpBinInt$_\vee$}{opbinint_vee}
$\vee : \mathbb{B} \times \mathbb{B} \to \mathbb{B}$ es una operación binaria interna.
\end{preaxioma}

\begin{preaxioma}{Operación Binaria Interna $\wedge$ - OpBinInt$_\wedge$}{opbinint_wedge}
$\wedge : \mathbb{B} \times \mathbb{B} \to \mathbb{B}$ es una operación binaria interna.
\end{preaxioma}

Para poder usar estos conceptos con mayor seguridad y flexibilidad en las futuras demostraciones formales, asignaremos nombres cortos a las condiciones de existencia y unicidad de la imagen para estas operaciones:

\begin{preaxioma}{Existencia $\vee$ - Existencia$_\vee$}{exist_vee}
Para todo par existe imagen en $\mathbb{B}$. Es decir, $\forall \langle a,b \rangle \in \mathbb{B} \times \mathbb{B}$, $\exists c \in \mathbb{B}$ tal que $a \vee b = c$.
\end{preaxioma}

\begin{preaxioma}{Existencia $\wedge$ - Existencia$_\wedge$}{exist_wedge}
Análogamente, $\forall \langle a,b \rangle \in \mathbb{B} \times \mathbb{B}$, $\exists d \in \mathbb{B}$ tal que $a \wedge b = d$.
\end{preaxioma}

\begin{preaxioma}{Unicidad $\vee$ - Unicidad$_\vee$}{unic_vee}
Para un par solo existe una imagen. Esto es, si $a \vee b = c$ y $a \vee b = d$, entonces $c = d$.
\end{preaxioma}

\begin{preaxioma}{Unicidad $\wedge$ - Unicidad$_\wedge$}{unic_wedge}
De igual forma para el ínfimo, si $a \wedge b = c$ y $a \wedge b = d$, entonces $c = d$.
\end{preaxioma}

\section{Los Postulados de Huntington (1904)}
Sobre el sistema $(\mathbb{B}, \vee, \wedge, \bot, \top)$ que cumple los pre-axiomas anteriores, diremos que forma un álgebra de Boole si satisface los siguientes postulados:

\begin{postulado}{Elemento neutro $\vee$ - $ElemNeu_\vee$}{neutro_vee}
Todo elemento operado mediante $\vee$ con el mínimo $\bot$ da como resultado el mismo elemento; es decir, $\bot$ no altera el valor original:
\[ \forall a \in \mathbb{B}, \quad a \vee \bot = a \]
\end{postulado}

\begin{postulado}{Elemento neutro $\wedge$ - $ElemNeu_\wedge$}{neutro_wedge}
Todo elemento operado mediante $\wedge$ con el máximo $\top$ da como resultado el mismo elemento, quedando inalterado:
\[ \forall a \in \mathbb{B}, \quad a \wedge \top = a \]
\end{postulado}

\begin{postulado}{Conmutatividad $\vee$ - $Comm_\vee$}{conmut_vee}
El orden de los operandos al aplicar la operación $\vee$ es indiferente, obteniéndose exactamente el mismo resultado:
\[ \forall a, b \in \mathbb{B}, \quad a \vee b = b \vee a \]
\end{postulado}

\begin{postulado}{Conmutatividad $\wedge$ - $Comm_\wedge$}{conmut_wedge}
De la misma forma, el orden de los operandos al aplicar la operación $\wedge$ tampoco altera el resultado final:
\[ \forall a, b \in \mathbb{B}, \quad a \wedge b = b \wedge a \]
\end{postulado}

\begin{postulado}{Distributividad $\vee$ sobre $\wedge$ - $Dist_\vee$}{distrib_vee_wedge}
La operación $\vee$ se distribuye sobre la operación $\wedge$. Operar un elemento con el resultado de un $\wedge$ equivale a operar con $\vee$ cada componente individualmente y luego aplicar $\wedge$:
\[ \forall a, b, c \in \mathbb{B}, \quad a \vee (b \wedge c) = (a \vee b) \wedge (a \vee c) \]
\end{postulado}

\begin{postulado}{Distributividad $\wedge$ sobre $\vee$ - $Dist_\wedge$}{distrib_wedge_vee}
De manera equivalente, el ínfimo ($\wedge$) se reparte de forma distributiva entre los componentes de un supremo ($\vee$):
\[ \forall a, b, c \in \mathbb{B}, \quad a \wedge (b \vee c) = (a \wedge b) \vee (a \wedge c) \]
\end{postulado}

\begin{postulado}{Complementario - $Comp_\vee, Comp_\wedge$}{comp}
Todo elemento del conjunto posee al menos un "complemento" (o elemento opuesto). Al operarlo con su complemento mediante $\vee$ siempre alcanzamos el máximo $\top$, y mediante $\wedge$ siempre caemos al mínimo $\bot$:
\begin{align*}
\forall a \in \mathbb{B}, \exists b \in \mathbb{B} \quad : \quad a \vee b &= \top \quad (Comp_\vee) \\
a \wedge b &= \bot \quad (Comp_\wedge)
\end{align*}
\end{postulado}

\textit{Nota: A diferencia de algunas formulaciones clásicas que imponen un axioma de cardinalidad ($\bot \neq \top$) para evitar el álgebra trivial, en este desarrollo permitiremos la existencia del álgebra trivial.}

